% Redrawn 2026-10-15; main labels 8 pt, notes 7.5 pt; actual fonts embedded in PDF.
% export-reference-width-mm: 169
\begin{tikzpicture}[font=\figurefont\fontsize{8}{10}\selectfont,
  box/.style={line width=.6pt,rounded corners=1.8pt,draw=FigureStroke,fill=FigureBlueFill,align=center,
    text width=66mm,minimum height=14mm},
  arr/.style={-{Latex[length=1.8mm]},draw=FigureStroke,line width=.6pt}]
  \begin{scope}[x=9mm,y=9mm]
    \fill[FigureYellowFill] (3.5,1.5) rectangle (4.5,3.5);
    \fill[BookRule!50] (1.5,1.5) rectangle (2.5,2.5);
    \foreach \x in {.5,1.5,2.5,3.5,4.5} \draw[BookRule] (\x,.5)--(\x,3.5);
    \foreach \y in {.5,1.5,2.5,3.5} \draw[BookRule] (.5,\y)--(4.5,\y);
    \foreach \x/\y/\n in {1/1/40,2/1/24,3/1/12,4/1/4,1/2/7,3/2/2,4/2/0,1/3/5,2/3/3,3/3/1,4/3/0}
      \node at (\x,\y) {\n};
    \node[below,align=center] at (2.5,0)
      {人工状态计数\\$N(s)=\sum_aN(s,a)$};
    \node[above] at (2.5,4) {(a) 真实访问覆盖};
  \end{scope}
  \node[box,fill=FigureGreenFill] (next) at (10,2.7)
    {后继乐观值：\\$V^+(4,1)=0.9$，$V^+(3,2)=0.4$\\$V^+(3,1)=0.2$};
  \node[box,fill=FigureYellowFill] (root) at (10,.1)
    {$(3,1)$向右的教学备份\\$-0.04+0.9(0.8\cdot0.9$\\$+0.1\cdot0.4+0.1\cdot0.2)+0.1$\\$=0.762$};
  \draw[arr] (next) -- node[right] {传播后继价值，再加置信奖励} (root);
  \node at (10,4) {(b) 乐观性沿合法转移传播};
  \node[align=center,text width=142mm] at (6.7,-2)
    {少访问不等于价值低；乐观上界帮助决定去哪里取得新证据。\\右侧数值仅说明一次备份，不是网格最优解或实测结果。};
\end{tikzpicture}
